\documentclass[12pt]{article}
\usepackage{fontspec}
\setmainfont{Charis SIL}
\usepackage[letterpaper, margin=1in]{geometry}
\usepackage{hyperref}
\usepackage{csquotes}

\begin{document}
\thispagestyle{empty}

\noindent\today

\bigskip

\noindent Dear Editors,

\medskip

I am submitting \enquote{Grammaticality needs more than dependency locality} for consideration as a Letter to the Editor in \textit{Trends in Cognitive Sciences}, in response to Gibson's article \enquote{Dependency syntax as the simplest theory of grammar} (\textit{TiCS} 30:4, 2026).

Gibson makes two claims. The first is representational: a language's syntax consists of dependency rules, with no phrase structure or transformations. The second is explanatory: dependency locality is the unifying principle behind processing effects and word-order universals. The representational case and the locality evidence are both strong, and I do not contest them. The Letter argues that the broader title claim, that dependency grammar is the simplest theory of grammar, overreaches. Three dissociations press against a single-factor story: the English string \textit{I have 57 years} carries the same dependency skeleton on a duration reading (licensed) and an age reading (not licensed in contemporary Standard English); six of seven English islands satiate under repeated exposure in Lu, Frank \& Degen's (2024) meta-analysis while the Left Branch Condition does not; and agreement-attraction effects concentrate in ungrammatical strings in ways a locality-only account does not predict. Gibson's own concluding remarks draw a narrower conclusion, describing dependency grammar as the simplest formalism for how words are connected within a sentence. That narrower claim, I argue, is the one the evidence supports.

The manuscript runs 781 words in the body, under the 800-word Letter cap, with 8 references, well under the 12-reference cap. There is no abstract, no figures or tables, and no section headings in the body, in keeping with the Letters format. Keywords: dependency grammar; dependency locality; grammaticality; island satiation; agreement attraction.

This manuscript is original, has not been published previously, and is not under consideration elsewhere. A preprint version has been posted on LingBuzz (\url{https://lingbuzz.net/lingbuzz/009938}) under a CC BY 4.0 license; the title has since been tightened to the current version.

As sole author I have no competing interests to declare, and the work received no external funding. The paper is theoretical and reports no new empirical data. Per Elsevier policy, the Letter's AI-assistance declaration names the tools used (Claude, ChatGPT) and the scope of use (structural review, critical feedback, and prose-editing suggestions), with full responsibility for the content resting with me.

Thank you for your consideration.

\bigskip

\noindent Sincerely,\\
Brett Reynolds\\
Humber Polytechnic \& University of Toronto\\
ORCID: 0000-0003-0073-7195\\
\href{mailto:brett.reynolds@humber.ca}{brett.reynolds@humber.ca}

\end{document}
